The systematic measurement of river stage and discharge is essential to having a baseline understanding of hydrological norms and extremes. They also
provide the critical information necessary for simulation models to be compared, calibrated, and improved. These models, both local, national, and
global serve as the foundation for flood forecasting and early warning systems, drought monitoring, water supply allocation, infrastructure design,
hydropower operations, and environmental flow assessment. Without reliable gauge records, the accuracy and credibility of both operational forecasts
and scientific analyses of river behavior are fundamentally limited.

Despite this importance, the global distribution of stream gauges is uneven, and the number of stations actively contributing to archives has been
declining since the 1980s \citep{ruhi2018}. Established in 1988, the Global Runoff Data Centre (GRDC) is the most prominent effort to compile
river discharge observations at the global scale \citep{grdc2024}. The GRDC currently holds daily and monthly discharge records from more than 9,800
stations across 160 countries, representing approximately 435,000 station-years of data with an average record length of 44 years. However, even the
GRDC's impressive data repository has only about 14\% of stations remaining active as of 2020, and coverage is heavily skewed toward North America
and Western Europe. While these regions maintain dense, well-funded monitoring networks, many developing regions, often those most vulnerable to
floods, droughts, and water scarcity, have inadequate coverage \citep{mishra2009}. Furthermore, data-sharing barriers related to national policy,
transboundary disputes, and institutional fragmentation mean that even where gauges exist, their records may not be accessible to the broader
scientific community or to global modeling efforts.

These gaps pose a challenge for large-scale hydrologic models such as the GEOGLOWS ECMWF River Forecast System, whose accuracy depends on observed
data for bias correction, calibration, and validation \citep{hales2022,sanchezlozano2021}
. Building a comprehensive, well-documented global gauge dataset is needed for not only improving global streamflow predictions, but also for
advancing our collective ability to manage water resources, reduce disaster risk, and improve water security in data-scarce regions worldwide.

The two purposes of this paper are first to describe a dataset of in situ river stage and discharge measurements compiled for global hydrology
modeling, and second to analyze the data for its strengths and weaknesses and correlations with socioeconomic indicators. The gauge measurements
themselves are not all publicly available, so we do not publish the full dataset for download in connection with this article. We instead provide a
collection of scripts and documentation to retrieve the data described separately. We describe the dataset in terms of its various summary properties
including its spatial coverage in terms of political boundaries as well as watersheds, and also the temporal availability of the measurements. The
detailed descriptors provide a basis for understanding the dataset without needing to distribute the entire catalog of measurements.
